% Generated from lessons/0011-3-continuous-function-properties.html; edit the HTML source, then rerun the converter.
\section{连续函数的性质}
\lessonmark{连续函数的性质}

这一节把一元闭区间上连续函数的性质推广到 \(R^n\)：紧集上连续映射保持紧性，实值连续函数有界并取到最大最小值，紧集上连续函数一致连续，连通集的连续像仍连通。

\subsection*{本节一句话}

\begin{conceptbox}
\textbf{紧集 + 连续}
连续映射把紧集映成紧集。实值时，紧集上的连续函数一定有界并取到最大最小值。
\end{conceptbox}

\begin{conceptbox}
\textbf{紧集 + 连续}
还会自动得到一致连续。局部的 \(\delta\) 可以通过有限覆盖变成全局统一的 \(\delta\)。
\end{conceptbox}

\begin{conceptbox}
\textbf{连通集 + 连续}
连续映射把连通集映成连通集。实值时，连通紧集上的值域是一个闭区间。
\end{conceptbox}

\begin{notebox}
本节考试关键词：紧集、连续、有界、最值、一致连续、连通、中间值。看到“紧集”先想“有限子覆盖/收敛子列”；看到“连通”先想“路径和中间值”。
\end{notebox}

\subsection*{预备：紧集性质速记}

在 \(R^n\) 中，紧集可以理解成高维空间里的闭区间。最重要的判定是 Heine-Borel 定理：

\[
      K\subset R^n\text{ 紧}
      \Longleftrightarrow
      K\text{ 闭且有界}.
    \]

\begin{center}
\small
\begin{tabularx}{\linewidth}{|>{\raggedright\arraybackslash}p{0.18\linewidth}|>{\raggedright\arraybackslash}p{0.42\linewidth}|>{\raggedright\arraybackslash}X|}
\hline
\textbf{性质} & \textbf{含义} & \textbf{考试用法} \\ \hline
闭性 & 紧集包含自己的极限点。 & 若 \(x_k\in K,\ x_k\to x\)，则 \(x\in K\)。 \\ \hline
有界性 & 紧集不能无限跑远。 & 可以把问题限制在某个大球内。 \\ \hline
序列紧性 & \(K\) 中任意点列都有收敛子列，且极限仍在 \(K\) 中。 & 证明连续像紧、最值存在时常用。 \\ \hline
聚点性质 & 紧集中的任意无限子集，在 \(K\) 中至少有一个聚点。 & 处理无限点集、反证有界闭性时常用。 \\ \hline
\end{tabularx}
\end{center}

\[
      x_k\in K
      \Longrightarrow
      \exists x_{k_j}\to x\in K.
    \]

\subsubsection*{和 Bolzano-Weierstrass 定理的关系}

Bolzano-Weierstrass 定理说：\(R^n\) 中任意有界点列，都存在收敛子列。它只负责“抓住一个收敛子列”，但不保证极限还在原集合里。

紧集的序列性质更强：点列不仅有收敛子列，而且这个子列的极限必须留在 \(K\) 里。

\[
      \text{紧集序列性质}
      =
      \text{Bolzano-Weierstrass 定理}
      +
      \text{闭性}.
    \]

例如 \(x_k=1/k\in(0,1)\)，它有收敛子列并且极限为 \(0\)，但 \(0\notin(0,1)\)。所以 \((0,1)\) 不是紧集。

\begin{notebox}
一句话记忆：Bolzano-Weierstrass 负责“不跑散”，闭性负责“极限不漏出去”。紧集同时拥有这两件事。
\end{notebox}

\subsection*{一、集合上的连续映射}

教材先把“在开集内连续”推广到任意点集 \(K\) 上。设 \(K\subset R^n\)，\(f:K\to R^m\)，\(x_0\in K\)。若对任意 \(\varepsilon>0\)，存在 \(\delta>0\)，使得

\[
      x\in O(x_0,\delta)\cap K
      \Longrightarrow
      |f(x)-f(x_0)|<\varepsilon,
    \]

则称 \(f\) 在 \(x_0\) 连续。若在 \(K\) 的每一点连续，则称 \(f\) 是 \(K\) 上的连续映射。

\begin{notebox}
这里的关键是 \(O(x_0,\delta)\cap K\)。如果 \(x_0\) 是 \(K\) 的边界点，只要求从 \(K\) 内部靠近 \(x_0\)，不要求从 \(K\) 外部靠近。
\end{notebox}

\subsubsection*{序列判别法}

在 \(R^n\) 中，连续性可以用序列来判别：

\[
      f\text{ 在 }a\in K\text{ 连续}
      \Longleftrightarrow
      \forall x_k\in K,\ x_k\to a,\quad f(x_k)\to f(a).
    \]

这是第 1 题的核心，也是后面证明“连续像保持紧性”的工具。

\subsection*{二、紧集上的连续映射}

\subsubsection*{定理 11.3.1：连续像保持紧性}

若 \(K\subset R^n\) 是紧集，\(f:K\to R^m\) 是连续映射，则

\[
      f(K)=\{y\in R^m:y=f(x),\ x\in K\}
    \]

也是紧集。

\subsubsection*{证明思路}

\begin{enumerate}
 \item 任取 \(f(K)\) 中的点列 \(\{y_k\}\)。
 \item 每个 \(y_k\) 都来自某个 \(x_k\in K\)，即 \(f(x_k)=y_k\)。
 \item 因为 \(K\) 紧，\(\{x_k\}\) 有收敛子列 \(x_{k_j}\to a\in K\)。
 \item 由连续性，\(y_{k_j}=f(x_{k_j})\to f(a)\in f(K)\)。
 \item 所以 \(f(K)\) 中任意点列都有收敛子列且极限仍在 \(f(K)\)，故 \(f(K)\) 紧。
 \end{enumerate}

\subsubsection*{定理 11.3.2 与 11.3.3：有界性与最值}

若 \(K\subset R^n\) 是紧集，\(f:K\to R\) 是连续函数，则 \(f\) 在 \(K\) 上有界；并且存在 \(\xi_1,\xi_2\in K\)，使得对任意 \(x\in K\)，

\[
      f(\xi_1)\le f(x)\le f(\xi_2).
    \]

也就是说，紧集上的实值连续函数一定取到最小值和最大值。

\begin{notebox}
一元闭区间上的最大最小值定理，在这里变成“紧集上的最大最小值定理”。在 \(R^n\) 中，“紧”可以理解成“闭且有界”。
\end{notebox}

\subsection*{三、一致连续}

设 \(K\subset R^n\)，\(f:K\to R^m\)。如果对任意 \(\varepsilon>0\)，存在 \(\delta>0\)，使得对 \(K\) 中任意两点 \(x',x''\)，只要

\[
      |x'-x''|<\delta,
    \]

就有

\[
      |f(x')-f(x'')|<\varepsilon,
    \]

则称 \(f\) 在 \(K\) 上一致连续。

\begin{center}
\small
\begin{tabularx}{\linewidth}{|>{\raggedright\arraybackslash}p{0.18\linewidth}|>{\raggedright\arraybackslash}p{0.42\linewidth}|>{\raggedright\arraybackslash}X|}
\hline
\textbf{性质} & \textbf{普通连续} & \textbf{一致连续} \\ \hline
\(\delta\) 是否依赖点 & 可以依赖 \(x_0\)。 & 不能依赖点，对整个 \(K\) 共用一个 \(\delta\)。 \\ \hline
直观 & 每个点附近都不跳。 & 整个集合上变化速度都能统一控制。 \\ \hline
常用定理 & 局部性质。 & 紧集上的连续函数必一致连续。 \\ \hline
\end{tabularx}
\end{center}

\subsubsection*{定理 11.3.4：紧集上一致连续性定理}

若 \(K\subset R^n\) 是紧集，\(f:K\to R^m\) 是连续映射，则 \(f\) 在 \(K\) 上一致连续。

\subsubsection*{证明思路}

\begin{enumerate}
 \item 对每个 \(a\in K\)，由连续性选一个局部半径 \(\delta_a\)，使得 \(x\) 靠近 \(a\) 时 \(|f(x)-f(a)|<\varepsilon/2\)。
 \item 这些小球 \(O(a,\delta_a/2)\) 覆盖 \(K\)。
 \item 因为 \(K\) 紧，可以取有限个球覆盖 \(K\)。
 \item 取这些有限半径的最小值的一半作为统一 \(\delta\)。
 \item 若 \(x',x''\in K\) 且 \(|x'-x''|<\delta\)，二者落入同一个控制球附近，用三角不等式得到 \(|f(x')-f(x'')|<\varepsilon\)。
 \end{enumerate}

\begin{notebox}
证明的一句话骨架：点点连续给出许多个局部 \(\delta_a\)，紧性把无限多个局部控制压缩成有限个，再取最小值，得到统一 \(\delta\)。
\end{notebox}

\subsection*{四、连通集与中间值}

教材这里用“道路”来定义连通。若存在连续映射

\[
      \gamma:[0,1]\to R^n,\qquad \gamma([0,1])\subset S,
    \]

且 \(\gamma(0)=x,\ \gamma(1)=y\)，则称 \(\gamma\) 是 \(S\) 中从 \(x\) 到 \(y\) 的道路。若 \(S\) 中任意两点之间都有道路，则称 \(S\) 为连通集。

连通的开集称为区域；区域的闭包称为闭区域。

\subsubsection*{定理 11.3.5：连续像保持连通性}

连续映射将连通集映成连通集。也就是说，若 \(D\subset R^n\) 连通，\(f:D\to R^m\) 连续，则 \(f(D)\) 连通。

\subsubsection*{证明思路}

取 \(f(D)\) 中任意两点 \(f(x),f(y)\)。由于 \(D\) 连通，存在道路 \(\gamma\) 从 \(x\) 连到 \(y\)。于是复合映射

\[
      f\circ\gamma:[0,1]\to R^m
    \]

就是 \(f(D)\) 中从 \(f(x)\) 连到 \(f(y)\) 的道路。因此 \(f(D)\) 连通。

\subsubsection*{定理 11.3.6：中间值定理}

若 \(K\subset R^n\) 是连通的紧集，\(f:K\to R\) 是连续函数，则 \(f\) 能取到最小值 \(m\) 与最大值 \(M\) 之间的一切值。换言之，

\[
      f(K)=[m,M].
    \]

\begin{notebox}
连通给“不中断”，紧集给“端点也取到”。所以连通紧集上的实值连续函数，值域不是散开的点集，而是一个闭区间。
\end{notebox}

\subsection*{五、题型模板}

\begin{conceptbox}
\textbf{证明有界或最值存在}
先找到一个紧集，再说明函数在这个紧集上连续。结论立刻是有界、取到最大最小值。
\end{conceptbox}

\begin{conceptbox}
\textbf{证明一致连续}
若定义域是紧集，直接用定理。若不是紧集，就要额外估计，例如 Lipschitz 条件。
\end{conceptbox}

\begin{conceptbox}
\textbf{证明集合开闭}
常把集合写成连续函数的原像：\(\{x:f(x)< c\}\) 是开集，\(\{x:f(x)\le c\}\) 是闭集。
\end{conceptbox}

\begin{conceptbox}
\textbf{证明距离函数性质}
用三角不等式证明 \(|d(x,A)-d(y,A)|\le |x-y|\)，这会直接给出一致连续。
\end{conceptbox}

\subsection*{六、快速检测}

\begin{quickcheck}
连续函数在什么类型的集合上一定取到最大值和最小值？

\tcblower
\textbf{答案：}紧集
\end{quickcheck}

\begin{quickcheck}
连续函数在任意开集上是否一定一致连续？

\tcblower
\textbf{答案：}不一定
\end{quickcheck}

\begin{quickcheck}
连续映射会把哪类集合映成同类集合，从而推出中间值性质？

\tcblower
\textbf{答案：}连通集
\end{quickcheck}

\subsection*{七、课后题训练}

下面按教材第 119 页习题 1-8 整理。数学式保留，提示用于复习证明主线；写作业时以教材原题为准。

\begin{exercisebox}
\textbf{习题 1\badge{连续的序列判别}}

设 \(D\subset R^n\)，\(f:D\to R^m\) 为连续映射。如果 \(D\) 中的点列 \(\{x_k\}\) 满足

\[
        \lim_{k\to\infty}x_k=a,\qquad a\in D,
      \]

证明：

\[
        \lim_{k\to\infty}f(x_k)=f(a).
      \]

\begin{hintbox}
直接用 \(f\) 在 \(a\) 连续的定义。给定 \(\varepsilon\)，先取 \(\delta\)，再由 \(x_k\to a\) 得到充分大的 \(k\) 使 \(|x_k-a|<\delta\)。
\end{hintbox}
\end{exercisebox}

\begin{exercisebox}
\textbf{习题 2\badge{原像开闭}}

设 \(f\) 是 \(R^n\) 上的连续函数，\(c\) 为实数。设

\[
        A_c=\{x\in R^n:f(x)< c\},\qquad
        B_c=\{x\in R^n:f(x)\le c\}.
      \]

证明：\(A_c\) 为 \(R^n\) 上的开集，\(B_c\) 为 \(R^n\) 上的闭集。

\begin{hintbox}
方法一：\(A_c=f^{-1}((-\infty,c))\)，开集的连续原像开；\(B_c=f^{-1}((-\infty,c])\)，闭集的连续原像闭。

方法二：对 \(A_c\) 中点 \(x_0\)，令 \(\varepsilon=c-f(x_0)>0\)。对 \(B_c\) 用序列判别闭集。
\end{hintbox}
\end{exercisebox}

\begin{exercisebox}
\textbf{习题 3\badge{连续但非一致连续}}

设二元函数

\[
        f(x,y)=\frac{1}{1-xy},\qquad (x,y)\in D=[0,1)\times[0,1).
      \]

证明：\(f\) 在 \(D\) 上连续，但不一致连续。

\begin{hintbox}
连续性来自分母 \(1-xy\ne0\)。不一致连续可取靠近 \((1,1)\) 的两列点，例如让两点距离趋于 \(0\)，但函数值差不趋于 \(0\)。这题的本质是：\(D\) 不是紧集，函数在边界附近爆掉。
\end{hintbox}
\end{exercisebox}

\begin{exercisebox}
\textbf{习题 4\badge{到集合的距离}}

设 \(A\) 为 \(R^n\) 上的非空子集，定义 \(R^n\) 上的函数 \(f\) 为

\[
        f(x)=\inf\{|x-y|:y\in A\},
      \]

它称为 \(x\) 到 \(A\) 的距离。证明：

\begin{enumerate}
 \item 当且仅当 \(x\in\overline A\) 时，\(f(x)=0\)。
 \item 对于任意 \(x',x''\in R^n\)，不等式 \[
            |f(x')-f(x'')|\le |x'-x''|
          \] 成立，从而 \(f\) 在 \(R^n\) 上一致连续。
 \item 若 \(A\) 是紧集，则对于任意 \(c>0\)，点集 \[
            \{x\in R^n:f(x)\le c\}
          \] 是紧集。
 \end{enumerate}

\begin{hintbox}
(1) 距离为 0 意味着可以找到 \(A\) 中点列趋于 \(x\)，这就是 \(x\in\overline A\)。

(2) 对任意 \(y\in A\)，由三角不等式 \(|x'-y|\le |x'-x''|+|x''-y|\)，再取下确界。

(3) 先由连续性得闭，再用 \(A\) 有界证明这个 \(c\)-邻域也有界。
\end{hintbox}
\end{exercisebox}

\begin{exercisebox}
\textbf{习题 5\badge{无穷远处极限与最值}}

设二元函数 \(f\) 在 \(R^2\) 上连续，证明：

\begin{enumerate}
 \item 若 \[
            \lim_{x^2+y^2\to+\infty}f(x,y)=+\infty,
          \] 则 \(f\) 在 \(R^2\) 上的最小值必定存在。
 \item 若 \[
            \lim_{x^2+y^2\to+\infty}f(x,y)=0,
          \] 则 \(f\) 在 \(R^2\) 上的最大值与最小值至少存在一个。
 \end{enumerate}

\begin{hintbox}
(1) 无穷远处趋于 \(+\infty\)，所以极小值只能出现在某个足够大的闭圆盘内；闭圆盘紧，连续函数取最小值。

(2) 若某点函数值 \(>0\)，最大值会被迫出现在大圆盘内；若没有正值但有负值，最小值会出现在大圆盘内；若恒为 0，则最大最小值都存在。
\end{hintbox}
\end{exercisebox}

\begin{exercisebox}
\textbf{习题 6\badge{齐次函数估计}}

设 \(f\) 是 \(R^n\) 上的连续函数，满足：

\begin{enumerate}
 \item 当 \(x\ne0\) 时成立 \(f(x)>0\)。
 \item 对于任意 \(x\) 与 \(c>0\)，成立 \(f(cx)=cf(x)\)。
 \end{enumerate}

证明：存在 \(a>0,\ b>0\)，使得

\[
        a|x|\le f(x)\le b|x|.
      \]

\begin{hintbox}
在单位球面 \(S=\{x:|x|=1\}\) 上使用最值定理。因为 \(S\) 紧且 \(f>0\)，可取到正的最小值 \(a\) 和最大值 \(b\)。对 \(x\ne0\)，写成 \(x=|x|\cdot x/|x|\)。
\end{hintbox}
\end{exercisebox}

\begin{exercisebox}
\textbf{习题 7\badge{闭包与连续像}}

设 \(f:R^n\to R^m\) 为连续映射。证明：对于 \(R^n\) 中的任意子集 \(A\)，成立

\[
        f(\overline A)\subset \overline{f(A)}.
      \]

举例说明 \(f(\overline A)\) 能够是 \(\overline{f(A)}\) 的真子集。

\begin{hintbox}
若 \(x\in\overline A\)，取 \(A\) 中点列 \(x_k\to x\)。连续性给出 \(f(x_k)\to f(x)\)，所以 \(f(x)\in\overline{f(A)}\)。

严格包含例子：取 \(A=R\)，\(f(x)=\arctan x\)。则 \(f(\overline A)=(-\pi/2,\pi/2)\)，而 \(\overline{f(A)}=[-\pi/2,\pi/2]\)。
\end{hintbox}
\end{exercisebox}

\begin{exercisebox}
\textbf{习题 8\badge{一致连续延拓}}

设 \(f\) 是有界开区域 \(D\subset R^n\) 上的一致连续函数，证明：

\begin{enumerate}
 \item 可以将 \(f\) 连续延拓到 \(D\) 的边界上，即存在定义在 \(\overline D\) 上的连续函数 \(\widetilde f\)，使得 \(\widetilde f|_D=f\)。
 \item \(f\) 在 \(D\) 上有界。
 \end{enumerate}

\begin{hintbox}
对边界点 \(\xi\)，取 \(D\) 中点列 \(x_k\to\xi\)。一致连续保证 \(f(x_k)\) 是 Cauchy 列，因而有极限。再证明这个极限与所取点列无关，用它定义 \(\widetilde f(\xi)\)。

延拓后 \(\widetilde f\) 在 \(\overline D\) 上连续。因为 \(D\) 有界，所以 \(\overline D\) 是紧集，于是 \(\widetilde f\) 有界，故 \(f\) 在 \(D\) 上有界。
\end{hintbox}
\end{exercisebox}

来源：陈纪修、于崇华、金路《数学分析》第三版下册，第十一章 §3，教材第 115-119 页。若某题需要完整书面证明，可以直接在浏览器中选中题目问我。

上一节：第十一章 §2：多元连续函数。下一节：第十二章 §1：偏导数与全微分。路线图：Euclid 空间章节路线图。
